% Abhakseite „Das kann ich“ – Zone nur im Gesamt (\mitzone)
%% TODO Ich-kann-Sätze: Platzhalter wie in den Aufgabendateien
\begin{abhakseite}
\ifdefined\mitzone
\abhakgruppe{Kennst du schon}
\abhak{1}{Kästchen im Raster abzählen und Strecken auf dem Karo abtragen}
\abhak{2}{Senkrechte und parallele Geraden mit dem Geodreieck zeichnen und erkennen, Abstand eines Punktes von einer Geraden messen}
\abhak{3}{Negative Zahlen an der Zahlengeraden ablesen und eintragen}
\abhak{4}{Vierecksarten und Dreiecksarten an ihren Eigenschaften erkennen und benennen (Quadrat, Rechteck, Raute, Parallelogramm, Drachenviereck, gleichschenkliges Trapez; gleichseitiges und gleichschenkliges Dreieck)}
\abhak{5}{Diagonalen eines Vierecks einzeichnen und benennen}
\abhak{6}{Winkel messen und zeichnen, rechten Winkel erkennen und markieren}
\abhak{7}{Strecken messen und gleich lange Strecken abtragen}
\abhak{8}{Kästchen im Raster abzählen und Strecken auf dem Karo abtragen – Fehler finden}
\abhak{9}{Kästchen im Raster abzählen und Strecken auf dem Karo abtragen – selbst rechnen}
\fi
\abhakgruppe{Einheit 1 · Koordinatensystem}
\abhak{10}{„Erst rechts, dann hoch“}
\abhak{11}{Punkte und Figuren im Koordinatensystem}
\abhak{12}{Achsen und Ursprung benennen, Einteilung ablesen}
\abhak{13}{Schreibweise P(x | y) lesen und schreiben (erste Zahl nach rechts, zweite nach oben)}
\abhak{14}{Streckenlänge parallel zu einer Achse abzählen}
\abhak{15}{Punkte und Figuren im Koordinatensystem – Fehler finden, Begründen, Darstellungswechsel, Anwendung}
\abhakgruppe{Einheit 2 · Achsensymmetrie und Spiegeln}
\abhak{16}{„Wie viele Achsen?“}
\abhak{17}{Symmetrieachsen bestimmen}
\abhak{18}{Spiegeln und Spiegelbild benennen}
\abhak{19}{prüfen, ob eine eingezeichnete Gerade Symmetrieachse ist (falten oder messen)}
\abhak{20}{Symmetrie der Vierecksarten im Haus der Vierecke zuordnen}
\abhak{21}{Symmetrieachsen bestimmen – Fehler finden, Begründen, Darstellungswechsel, Anwendung}
\abhakgruppe{Einheit 3 · Punktsymmetrie, Drehung, Verschiebung}
\abhak{22}{Drehen, Verschieben, Punktsymmetrie}
\abhak{23}{drehsymmetrische Figuren erkennen und den Drehwinkel angeben (Vorrat)}
\abhak{24}{Original und Bild vergleichen: Längen und Winkel bleiben gleich, kongruente Figuren}
\abhak{25}{Drehen, Verschieben, Punktsymmetrie – Fehler finden, Begründen, Darstellungswechsel, Anwendung}
\end{abhakseite}
